% concatenated from polystab_manuscript.tex
\documentclass[a4paper,12pt]{amsart}
\pagestyle{myheadings}
\usepackage{amsmath,amssymb,amsfonts,amsthm}
\usepackage{latexsym,mathrsfs}
\usepackage{graphicx}
\usepackage[all]{xy}
\input xypic
\usepackage{color}
%\usepackage[backref=false]{hyperref}
\usepackage[pagebackref]{hyperref}
\hypersetup{colorlinks=true,linkcolor=red,citecolor=blue}
\usepackage[OT2,T1]{fontenc}
\usepackage[utf8]{inputenc}
\usepackage{lmodern}
\usepackage{microtype}
\usepackage{enumerate}
\usepackage{tikz-cd}
\usepackage{bm}
\usepackage{ulem}
\usepackage{mathtools}
\usepackage{tensor}
%%%%%%%%%%%%%%%%%%%
%\usepackage[notref,notcite]{showkeys}
% Ref. numbers; TO BE REMOVED IN THE FINAL VERSION!
%%%%%%%%%%%%%%%%%%%
%%%%%%%%%%%%%%%%%%%%%%%%%%%%%%%%%%%%
% Environments:
\theoremstyle{plain}
\newtheorem{theorem}[subsection]{{\bf Theorem}}
\newtheorem{conjecture}[subsection]{{\bf conjecture}}
%\newtheorem*{theorem*}{{\bf Theorem}}
\newtheorem{corollary}[subsection]{{\bf Corollary}}
%\newtheorem*{corollary*}{{\bf Corollary}}
\newtheorem{proposition}[subsection]{{\bf Proposition}}
\newtheorem{lemma}[subsection]{{\bf Lemma}}
%\newtheorem{fact}[subsection]{{\bf Fact}}
%\newtheorem{question}[subsection]{{\bf Question}}
%\newtheorem{conjecture}[subsection]{{\bf Conjecture}}
\theoremstyle{definition}
\newtheorem*{definition}{{\bf Definition}}
\theoremstyle{remark}
\newtheorem{remark}[subsection]{{\it Remark}}
\newtheorem{example}[subsection]{{\it Example}}
\numberwithin{equation}{subsection}
%%%%%%%%%%%%%%%%%%%%%%%%%%%%%%%%%%%%
%%%%%%%%%%%%%%%%%%%%%%%%%%%%%%%%%%%%
% Custom commands:
\DeclareMathOperator{\im}{im}
\DeclareMathOperator{\Tor}{Tor}
\DeclareMathOperator{\Ext}{Ext}
\DeclareMathOperator{\HH}{H}
\DeclareMathOperator{\K}{K}
\DeclareMathOperator{\B}{B}
\DeclareMathOperator{\Aut}{Aut}
\DeclareMathOperator{\Out}{Out}
\DeclareMathOperator{\Inn}{Inn}
\DeclareMathOperator{\Hom}{Hom}
\DeclareMathOperator{\End}{End}
\DeclareMathOperator{\GL}{GL}
\DeclareMathOperator{\SL}{SL}
\DeclareMathOperator{\UT}{UT}
\DeclareMathOperator{\GF}{GF}
\DeclareMathOperator{\Der}{Der}
\DeclareMathOperator{\res}{res}
\DeclareMathOperator{\cor}{cor}
\DeclareMathOperator{\conj}{conj}
\DeclareMathOperator{\tra}{tra}
\DeclareMathOperator{\cyl}{cyl}
\DeclareMathOperator{\cone}{cone}
\DeclareMathOperator{\coker}{coker}
\DeclareMathOperator{\Sym}{Sym}
\DeclareMathOperator{\Fix}{Fix}
\DeclareMathOperator{\id}{id}
\DeclareMathOperator{\Sub}{Sub}
%%%%%%%%%%%%%%%%%%%%%%%%%%%%%%%%%%%%
\DeclareBoldMathCommand{\bbot}{\bot}
\DeclareMathOperator{\fuse}{\mathfrak{fuse}}
\DeclareSymbolFont{cyrletters}{OT2}{wncyr}{m}{n}
\DeclareMathSymbol{\Sha}{\mathalpha}{cyrletters}{"58}
%%%%%%%%%%%%%%%%%%%%%%%%%%%%%%%%%%%%
\DeclareMathOperator{\pwc}{pwc}
\DeclareMathOperator{\pwh}{pwh}
\DeclareMathOperator{\cl}{cl}
\DeclareMathOperator{\ad}{ad}
%%%%%%%%%%%%%%%%%%%%%%%%%%%%%%%%%%%%
\DeclareMathOperator{\QHom}{QHom}
\DeclareMathOperator{\MQHom}{MQHom}
\DeclareMathOperator{\QAff}{QAff}
\DeclareMathOperator{\Poly}{Poly}
\DeclareMathOperator{\QPoly}{QPoly}
\DeclareMathOperator{\MQPoly}{MQPoly}
\DeclareMathOperator{\Conj}{Conj}
\DeclareMathOperator{\Comm}{Comm}
\DeclareMathOperator{\LComm}{LComm}
\DeclareMathOperator{\RComm}{RComm}
\DeclareMathOperator{\Pol2}{Pol_2}
\DeclareMathOperator{\Pold}{Pol}
\DeclareMathOperator{\CAT}{CAT}
\DeclareMathOperator{\sgn}{sgn}
\DeclareMathOperator{\alt}{alt}
\DeclareMathOperator{\grph}{\mathfrak{gr}}
\DeclareMathOperator{\poly}{poly}
\DeclareMathOperator{\HK}{HK}
\DeclareMathOperator{\deff}{def}
\DeclareMathOperator{\Alt}{Alt}
%\DeclareMathOperator{\mypartial}{\mathfrak{d}}
\DeclareMathOperator{\mypartial}{\partial}
\DeclareMathOperator{\taopartial}{\partial}
\newcommand{\llangle}{\langle\!\langle}
\newcommand{\rrangle}{\rangle\!\rangle}
%%%%%%%%%%%%%%%%%%%%%%%%%%%%%%%%%%%%
\newcommand{\QZ}{\mathbb{Q}/\mathbb{Z}}
\newcommand{\Z}{\mathbb{Z}}
\newcommand{\N}{\mathbb{N}}
\newcommand{\Q}{\mathbb{Q}}
\newcommand{\R}{\mathbb{R}}
\newcommand{\C}{\mathbb{C}}
\newcommand{\cwedge}{\curlywedge}
\newcommand{\dwedge}{\tilde{\wedge}}
%%%%%%%%%%%%%%%%%%%%%%%%%%%%%%%%%%%%
\newcommand{\norm}[1]{\left\lVert #1\right\rVert}
\newcommand{\abs}[1]{\left\lvert #1\right\rvert}
%%%%%%%%%%%%%%%%%%%%%%%%%%%%%%%%%%%%
% Bar complex:
\newcommand{\Cbar}{\bar{C}}
% Differential for homogeneous complex:
\newcommand{\dif}{\delta}
% Differential for bar complex:
\newcommand{\dbar}{\bar{\delta}}
%%%%%%%%%%%%%%%%%%%%%%%%%%%%%%%%%%%%%%%%%%%%%%%%%%%%%%%%%%%%%%%%%%%%%%%%
% Indices:
% cp
\newcommand{\cp}{\mspace{-1mu}{}_{cp}}
% Bounded cp:
\newcommand{\bc}{\mspace{-1mu}{}_{bc}}
% Bounded:
\newcommand{\bb}{\mspace{-1mu}{}_{b}}
%%%%%%%%%%%%%%%%%%%%%%%%%%%%%%%%%%%%%%%%%%%%%%%%%%%%%%%%%%%%%%%%%%%%%%%%
% Colored comments
\newcommand{\comment}[1]{\textcolor{purple}{#1}}
\newcommand{\answer}[1]{\textcolor{gray}{#1}}
\newcommand{\ai}[1]{\textcolor{brown}{#1}}
%%%%%%%%%%%%%%%%%%%%%%%%%%%%%%%%%%%%%%%%%%%%%%%%%%%%%%%%%%%%%%%%%%%%%%%%
%%%%%%%%%%%%%%%%%%%%%%%%%%%%%%%%%%%%%%%%%%%%%%%%%%%%%%%%%%%%%%%%%%%%%%%%
\begin{document}
\title[Polynomial stability and soficity]{Polynomial permutation stability, soficity, and universal polynomial groups}
\author{Primo\v z Moravec}
% \address{{
% Faculty of  Mathematics and Physics, University of Ljubljana,
% and Institute of Mathematics, Physics and Mechanics,
% Slovenia}}
% \email{primoz.moravec@fmf.uni-lj.si}
\thanks{Faculty of  Mathematics and Physics, University of Ljubljana,
and Institute of Mathematics, Physics and Mechanics,
Slovenia, 
\texttt{primoz.moravec@fmf.uni-lj.si}}
\subjclass[2020]{Primary 20F69; Secondary 20F70, 20F05}
\keywords{Polynomial maps of groups, permutation stability, sofic groups.}
\thanks{ORCID: \url{https://orcid.org/0000-0001-8594-0699}.}
\thanks{The author acknowledges the financial support from the Slovenian Research and Innovation Agency (ARIS), research core funding No. P1-0222, and project No. J1-50001.}
\date{\today}
%%%%%%%%%%%%%%%%%%%%%%%%%%%%%%%%%%%%%%%%%%%%%%%%%%%%%%%%%%%%%%%%%%%%%
\begin{abstract}
\noindent
We introduce polynomial permutation stability, extending classical
permutation stability from homomorphisms to polynomial maps of groups.
The universal polynomial group $\Pold_s(G)$ provides a natural
framework for this theory, and we prove that polynomial stability of
degree $s$ is equivalent to permutation stability of $\Pold_s(G)$.
This yields, in particular, polynomial stability of all finite groups
in degree two.

We develop a corresponding theory of polynomial sofic
approximations and show that, for countable groups, it does not give a
new notion of soficity. Nevertheless, the associated universal
polynomial groups contain new structural information: in degree two
they embed into a wreath product, which allows us to characterize their
soficity in terms of the original group. We also formulate and study
weak polynomial stability through the weak stability theory of
universal polynomial groups.

For finite groups, we investigate the structure of higher-degree
universal polynomial groups. We obtain general obstructions to
amenability and prove largeness results for a family of finite groups
with perfect derived subgroup and abelianization of order two.
\end{abstract}
% {\color{blue}
% We introduce polynomial permutation stability for group-valued polynomial maps. A sequence of maps
% \[
% f_n:G\longrightarrow \operatorname{Sym}(m_n)
% \]
% is called an almost polynomial of degree at most \(s\) if it is asymptotically unital and all \((s+1)\)-fold differences at the identity converge to the identity in normalized Hamming distance. We prove that \(G\) is polynomial permutation stable of degree \(s\) if and only if the universal polynomial group \(\operatorname{Pol}_s(G)\) is permutation stable. In particular, every finite group is permutation stable of degree two.

% We also introduce polynomial sofic approximations and show that, for countable groups, their existence is equivalent to ordinary soficity. In degree two we construct an embedding
% \[
% \operatorname{Pol}_2(G)\hookrightarrow G^{\mathrm{ab}}\wr G,
% \]
% and deduce that \(\operatorname{Pol}_2(G)\) is sofic if and only if \(G\) is sofic. We further relate weak polynomial stability to weak permutation stability of \(\operatorname{Pol}_s(G)\).

% Finally, for finite \(G\), we investigate amenability of the universal groups. If \(3\) divides \(\lvert G^{\mathrm{ab}}\rvert\), then \(\operatorname{Pol}_s(G)\) is nonamenable for every \(s\geq3\). For \(G=N\rtimes C_2\), where \(N\) is a nontrivial finite perfect group, we obtain an explicit relative presentation of \(\operatorname{Pol}_s(G)\) and prove that these groups are large for all sufficiently large \(s\).
% }
% \end{abstract}
%%%%%%%%%%%%%%%%%%%%%%%%%%%%%%%%%%%%%%%%%%%%%%%%%%%%%%%%%%%%%%%%%%%%%%%
\maketitle

% \tableofcontents
%%%%%%%%%%%%%%%%%%%%%%%%%%%%%%%%%%%%%%%%%%%%%%%%%%%%%%%%%%%%%%%%%%%%%%%
\section{Introduction}
\label{s:intro}

\noindent
Polynomial maps between groups, defined by the vanishing of iterated
finite differences, provide a nonlinear extension of the theory of group
homomorphisms.  They were systematically studied by Leibman
\cite{Lei02} and appear naturally in higher-order Fourier analysis and in
the study of Gowers norms \cite{GT12,GTZ12,JT24}.
Jamneshan and Thom recently obtained an
explicit description of the universal group representing unital
quadratic maps and developed several structural and stability results for
polynomial maps on nonabelian groups \cite{JT24}.  They also construct, for every $s\geq1$, a universal group $\Pold_s(G)$
representing unital polynomial maps of degree at most $s$.

The purpose of this paper is to connect these universal polynomial groups
with permutation stability and sofic approximation.  Recall that a
countable group is permutation stable ($P$-stable for short) if every almost homomorphism into
finite symmetric groups, equipped with normalized Hamming distance, is
close to genuine homomorphisms.  This notion originates in the work of
Glebsky--Rivera \cite{GR09} and has since been developed through
ultraproducts, invariant random subgroups, and action traces; see, among
others, \cite{AP15,BLT19,AP24}.  We replace asymptotic multiplicativity by
asymptotic vanishing of higher differences and introduce the corresponding
notion of polynomial permutation stability of degree $s$ ($P_s$-stability for short). Our first result shows that this higher-degree stability problem is exactly
an ordinary stability problem for the universal polynomial group:

\begin{theorem}
    \label{thm:Pscharacterization}
    A group $G$ is $P_s$-stable if and only if $\Pold_s(G)$ is $P$-stable.
\end{theorem}

This observation makes the existing theory of permutation stability
available in the polynomial setting.  

A second theme concerns polynomial analogues of sofic approximations.  We
call a separating almost-polynomial sequence of degree at most $s$ a
polynomial sofic approximation of degree $\le s$.  At first sight this could
define a larger class than the class of sofic groups.  We show that no new
class arises:

\begin{theorem}
    \label{thm:soficcharacterization}
    A countable group $G$ admits a polynomial sofic approximation of degree $\le s$ if and only $G$ is sofic.
\end{theorem}

The main ingredient is
a degree-reduction argument: the restriction of a degree-$s$ polynomial
map to the derived subgroup has degree at most $s-1$.  Iteration, together
with the closure of soficity under extensions by amenable quotients,
reduces Theorem \ref{thm:soficcharacterization} to the usual degree-one case.
In degree two, the explicit structure theorem of Jamneshan--Thom gives a
more precise result, namely that $G$ is sofic if and only if $\Pold_2(G)$ is sofic. The reverse direction still holds in higher degrees: If $\Pold_s(G)$ is sofic, then $G$ is sofic (Corollary \ref{cor:soficG}). Hence this, together with recently proven existence of non-sofic groups \cite{Ope26}, provides a new source of non-sofic groups. On the other hand, we do not know whether or not the soficity of $G$ implies the soficity of $\Pold_s(G)$ when $s\ge 3$.

A result of Glebsky and Rivera \cite{GR09} shows that finite groups are $P$-stable. Since Jamneshan--Thom \cite{JT24} proved that
$\Pold_2(G)$ is finite whenever $G$ is finite, Theorem \ref{thm:Pscharacterization} immediately yields that every finite group is
$P_2$-stable. On the other hand, a recent result of Alekseev and Thom \cite{AT26} shows that the group $\Pold_3(C_3)$ is far from being amenable, hence $P_s$-stability of finite groups remains open for $s\ge 3$. We address this question by introducing weak polynomial permutation stability of finite groups.  Following
Arzhantseva--P\u{a}unescu \cite{AP15}, one tests both the defining
relations and the non-relations of the canonical presentation of
$\Pold_s(G)$.  We prove that weak $P_s$-stability of $G$ is equivalent to
weak permutation stability of $\Pold_s(G)$.  Consequently, when
$\Pold_s(G)$ is amenable, weak $P_s$-stability is equivalent to residual
finiteness of $\Pold_s(G)$.  We also record a polynomial consequence of
the action-trace criterion of Arzhantseva--P\u{a}unescu \cite{AP24}. These results raise question of when the structure of $\Pold_s(G)$ can be controlled in a sense. The final part of the paper is therefore devoted to studying $\Pold_s(G)$ for finite groups $G$. The computation of
$\Pold_3(C_3)$ by Alekseev--Thom \cite{AT26} shows that higher-degree
universal groups can be infinite and nonamenable even for a finite cyclic
source.  Functoriality therefore gives the general obstruction: If $G$ is finite and 
$3\mid |G^{\mathrm{ab}}|$, then 
$\Pold_s(G)$ is non-amenable for every $s\geq3$.
We complement this with a family whose abelianization is $C_2$.  If
$N$ is a nontrivial finite perfect group and
$G=N\rtimes_\alpha C_2$, where $\alpha^2=1$, we obtain an explicit
relative presentation of $\Pold_s(G)$.  Using a finite-index free subgroup
and Lackenby's theorem on adjoining high-powered relations
\cite{Lac07}, we prove that $\Pold_s(G)$ is large, and hence
non-amenable, for all sufficiently large $s$.  A $2$-deficiency estimate
based on \cite{BT11} gives an explicit sufficient bound.
% ${\color{blue}
%   We show that the defining cocycle of $\Pold_2(G)$
% becomes a coboundary after enlarging the abelian kernel, which yields an
% embedding
% \[
% \Pold_2(G)\hookrightarrow G^{\mathrm{ab}}\wr G.
% \tag{1.3}
% \]
% The soficity of restricted wreath products of countable sofic groups
% \cite{HS18}, together with (1.2), implies
% \[
% \Pold_2(G)\text{ is sofic}
% \quad\Longleftrightarrow\quad
% G\text{ is sofic}.
% \tag{1.4}
% \]



% The paper is organized as follows.  Section~\ref{s:prelim} recalls
% polynomial maps, universal polynomial groups, permutation stability, and
% soficity.  Section~\ref{s:almostpolynomials} develops the basic calculus
% of almost polynomials and proves the degree-reduction lemma.  In
% Section~\ref{s:polystability} we introduce $P_s$-stability and prove
% (1.1).  Section~\ref{s:polysoficapprox} establishes (1.2)--(1.4).
% Section~\ref{s:weakpolypermut} treats weak polynomial stability.  The
% final section studies nonamenability and largeness of $\Pold_s(G)$ for
% finite groups.
% }

%%%%%%%%%%%%%%%%%%%%%%%%%%%%%%%%%%%%%%%%%%%%%%%%%%%%%%%%%%%%%%%%%%%%%%%

\section{Preliminaries}
\label{s:prelim}

\subsection{Polynomial maps}
\label{ss:poly_prelim}

For groups $G$ and $H$, let $H^G$ be the set of all maps
$f:G\to H$. Given $g\in G$, we define the left difference operator $\mypartial_g$ on $H^G$ by the rule
$$(\mypartial_g f)(x)=f(gx)f(x)^{-1}.$$
Given $g_1,\ldots ,g_k\in G$, we use the notation
$$\mypartial_{g_1,\ldots,g_k}f=\mypartial_{g_1}\circ\cdots\circ\mypartial_{g_k}f.$$
According Leibman \cite{Lei02}, the map $f$ is said to be \emph{polynomial of degree at most $s$} if
$$(\mypartial_{g_1,\ldots ,g_{s+1}}f)(x)=1$$
for all $x,g_1,\ldots ,g_{s+1}\in G$. We remark here that Leibman \cite{Lei02} uses the right difference operator $(D_gf)(x)=f(x)^{-1}f(xg)$, but it is easy to see that the definitions are equivalent due to the identity
$$(D_gf)(x)=f(x)^{-1}\cdot (\mypartial_{xgx^{-1}}f)(x)\cdot f(x).$$
Also note that the left difference operator is used by Green, Tao and Ziegler \cite{GTZ12} in their (more general) treatment of polynomial maps, while Green and Tao \cite{GT12} use the right difference operator.

Observe that
$$(\mypartial_gf)(x)
=f(g_1x)f(1)^{-1}f(1)f(x)^{-1}
=(\mypartial_{gx}f)(1)\cdot(\mypartial_xf)(1)^{-1}$$ implies
\begin{align*}
    (\mypartial_{g_1,\ldots ,g_{s+1}}f)(x) &=
    (\mypartial_{g_1}(\mypartial_{g_2,\ldots,g_{s+1}}f))(x)\\
    &= (\mypartial_{g_1x}(\mypartial_{g_2,\ldots,g_{s+1}}f))(1)
    \cdot (\mypartial_x(\mypartial_{g_2,\ldots,g_{s+1}}f))(1)^{-1}\\
    &= (\mypartial_{g_1x,g_2,\ldots,g_{s+1}}f)(1)
    \cdot (\mypartial_{x,g_2,\ldots,g_{s+1}}f)(1)^{-1}.
\end{align*}
This shows that $f$ is polynomial of degree $\le s$ if and only if 
$$(\mypartial_{g_1,\ldots ,g_{s+1}}f)(1)=1$$
for all $g_1,\ldots ,g_{s+1}\in G$. This aligns with the definition of polynomial maps used by Jamneshan and Thom \cite{JT24}. 
%%%%%%%%%%%%%%%%%%%%%%%%%%%%%%%

\subsection{Universal polynomial construction}
\label{ss:universal_prelim}

The material of this subsection is adapted from \cite{JT24}. Let $G$ be a group. Let $F_G$ be the free group on formal generators $\{ x_g\mid g\in G\setminus\{1\}\}$, and set $x_1=1$ by notation. Define a map $\iota :G\to F_G$ by $\iota(g)=x_g$.
Let $R_s(G)$ be the normal closure in $F_G$ of the set of all words $(\mypartial_{g_1,\ldots ,g_{s+1}}\iota)(1)$, where $g_1,\ldots ,g_{s+1}$ run through $G$. Put
$\Pold_s(G)=F_G/R_s(G)$.

\begin{proposition}[\cite{JT24}, Theorem 1.2, Remark 2.5]
    \label{prop:Pold}
    There exists a universal unital polynomial map $\iota_s:G\to\Pold_s(G)$ of degree $s$ such that, 
    for every unital polynomial map $f:G\to H$ of degree $\le s$, there exists a unique homomorphism $\psi_f: \Pold_s(G)\to H$ such that $f=\psi_f\circ \iota_s$.
\end{proposition}

Explicitly, if $q:F_G\to\Pold_s(G)$ is the quotient homomorphism, then $\iota_s=q\circ\iota:G\to\Pold_s(G)$ is the universal polynomial map mentioned in Proposition \ref{prop:Pold}. 
Note that $\iota_s$ is injective. Namely, the identity map $\id_G:G\to G$ is a unital polynomial of degree $\le s$, therefore $\psi_{\id_G}\circ\iota_s=\id_G$.
%%%%%%%%%%%%%%%%%%%%%%%%%%%%%%%

\subsection{Almost-homomorphisms}
\label{ss:almosthom_prelim}
 
This is mainly taken from Becker, Lubotzky, Thom \cite{BLT19}.
% see also Bradford \cite{Bra26}.

Equip the symmetric groups $\Sym(n)$ with the normalized Hamming distance:
$$d_n(\sigma,\tau)=1-\frac{1}{n}|\Fix(\sigma^{-1}\tau)|.$$
The metric $d_n$ is bi-invariant.

Let $G$ be a countable group. 
An \emph{almost-homomorphism} of $G$ is a sequence of maps 
$f_n:G\to\Sym (m_n)$ satisfying
$$\lim_{n\to\infty} d_{m_n}(f_n(g)f_n(h),f_n(gh))=0$$
for all $g,h\in G$. Note that almost-homomorphisms are \emph{asymptotically unital}, that is,
$$\lim_{n_to\infty} d_{m_n}(f_n(1),1)=0.$$

Let $(f_n)_n$ be an almost-homomorphism of $G$. Choose a nonprincipal ultrafilter $\mathcal{U}$ on $\mathbb{N}$, and form the metric ultraproduct \cite{Gol22}
$$\mathfrak{S}_{\mathcal{U}}=\prod_{\mathcal{U}} (\Sym(m_n),d_{m_n})=\frac{\prod_n \Sym(m_n)}{\left\{(\sigma_n)_n\mid \lim_{n\to\mathcal{U}}d_{m_n}(\sigma_n,1)=0\right\}}.$$
Define $f:G\to\mathfrak{S}_{\mathcal{U}}$ by
$$f(g)=[f_n(g)]_{\mathcal{U}}.$$
The fact that $(f_n)_n$ is an almost-homomorphism implies that $f$ is a homomorphism of groups. 

Conversely, 
let $G$ be a countable group and let
$\Theta:G\to\mathfrak{S}_\mathcal{U}$
be a homomorphism into a metric ultraproduct.  For each $g\in G$,
choose representatives
$f_n(g)\in\Sym(m_n)$
such that
$$\Theta(g)=[(f_n(g))_n]_{\mathcal U}.$$
Since $\Theta$ is a homomorphism, for every $g,h\in G$ one has
$$
\lim_{n\to\mathcal U}
d_{m_n}\bigl(f_n(gh),f_n(g)f_n(h)\bigr)=0.
$$
As $G$ is countable, a standard diagonal argument yields a subsequence $(n_k)$ such that
$$\lim_{k\to\infty}d_{m_{n_k}}\bigl(
f_{n_k}(gh),f_{n_k}(g)f_{n_k}(h)
\bigr) =0
$$
for every $g,h\in G$.  Thus, after passing to a subsequence, every
homomorphism of a countable group into a metric ultraproduct of
symmetric groups is represented by an almost homomorphism in the usual
sequential sense.  Moreover, any countable collection of additional
ultralimit conditions may be preserved in the same diagonalization. In
particular, if
$d_{\mathcal U}\bigl(\Theta(g),1\bigr)=1$ for all $g\neq 1$,
then the subsequence may be chosen so that
$$\lim_{k\to\infty}d_{m_{n_k}}\bigl(f_{n_k}(g),1\bigr) =1$$
for every $g\neq 1$.

%%%%%%%%%%%%%%%%%%%%%%%%%%%%%%%

\subsection{Permutation stability}
\label{ss:pstab_prelim}

The almost-homomorphism $(f_n)_n$ of $G$ is \emph{close to a homomorphism} if there exists a sequence of group homomorphisms $\rho_n:G\to \Sym(m_n)$ such that
$$\lim_{n\to \infty} d_{m_n}(\rho_n(g),f_n(g))=0$$
for all $g\in G$. A group $G$ is \emph{permutation stable} (\emph{$P$-stable} for short) if every almost-homomorphism of $G$ is close to a homomorphism. Equivalently, $G$ is $P$-stable if and only if every homomorphism $f:G\to\mathfrak{S}_\mathcal{U}$ is liftable, i.e., there exists a homomorphism $\phi:G\to\prod_n\Sym(m_n)$ such that $f=Q\circ \phi$, where $Q:\prod_n\Sym(m_n)\to\mathfrak{S}_\mathcal{U}$ is the canonical projection, see \cite[Theorem 4.2]{AP15}.

The following is well-known, we include a proof for completeness:

\begin{lemma}
\label{lem:Pstabrel}
Let $G=F/N$, $N=\langle\!\langle R\rangle\!\rangle$,
where $F$ is a free group  and $R\subseteq F$ is an arbitrary set of relators. Suppose that $G$ is $P$-stable.
Let $f_n\colon F\to\Sym(m_n)$
be homomorphisms such that
$$d_{m_n}\bigl(f_n(r),1\bigr)\longrightarrow 0$$
for every $r\in R$.
Then there exist homomorphisms
$\rho_n\colon G\to\Sym(m_n)$
such that
$$d_{m_n}\bigl(f_n(w),\rho_n(wN)\bigr)\longrightarrow 0$$
for every $w\in F$.
Equivalently, if $\pi\colon F\to G$ is the quotient map, then
$$d_{m_n}\bigl(f_n(w),(\rho_n\circ\pi)(w)\bigr)\longrightarrow0$$
for every $w\in F$.
\end{lemma}

\begin{proof}
Choose a set-theoretic section
$s\colon G\to F$
of the quotient map $\pi\colon F\to G$, and define
$$\varphi_n(g)=f_n(s(g)).$$
Fix $g,h\in G$. Since
$s(g)s(h)s(gh)^{-1}\in N=\langle\!\langle R\rangle\!\rangle$,
there exist $r_1,\ldots,r_k\in R$, $\varepsilon_i\in\{-1,1\}$, and
$u_1,\ldots,u_k\in F$ such that
$$s(g)s(h)s(gh)^{-1}=
\prod_{i=1}^k u_i r_i^{\varepsilon_i}u_i^{-1}.$$
Using that $f_n$ is a homomorphism and that the Hamming metric is
bi-invariant, we obtain
$$
\begin{aligned}
d_{m_n}\bigl(\varphi_n(g)\varphi_n(h),\varphi_n(gh)\bigr)
&=
d_{m_n}\bigl(f_n(s(g)s(h)s(gh)^{-1}),1\bigr)\\
&\leq
\sum_{i=1}^k d_{m_n}\bigl(f_n(r_i),1\bigr)
\longrightarrow0.
\end{aligned}
$$
Hence $(\varphi_n)_n$ is an almost-homomorphism of $G$.
Since $G$ is $P$-stable, there exist homomorphisms
$$\rho_n\colon G\to\Sym(m_n)$$
such that
$$d_{m_n}\bigl(\varphi_n(g),\rho_n(g)\bigr)\longrightarrow0$$
for every $g\in G$.

Now fix $w\in F$. Since
$w\,s(\pi(w))^{-1}\in N$,
the same argument gives
$$d_{m_n}\bigl(f_n(w),f_n(s(\pi(w)))\bigr)\longrightarrow0.$$
Therefore
$$
\begin{aligned}
d_{m_n}\bigl(f_n(w),\rho_n(\pi(w))\bigr)
&\leq
d_{m_n}\bigl(f_n(w),f_n(s(\pi(w)))\bigr)\\
&\quad+
d_{m_n}\bigl(\varphi_n(\pi(w)),\rho_n(\pi(w))\bigr)
\longrightarrow0.
\end{aligned}
$$
This proves the claim.
\end{proof}

% \begin{definition}[Partial homomorphism]
    
% \end{definition}

% \begin{definition}[Locally stable]
    
% \end{definition}
%%%%%%%%%%%%%%%%%%%%%%%%%%%%%%%

\subsection{Sofic groups}
\label{ss:sofic_prelim}

Let $G$ be a countable group. A \emph{sofic approximation} of $G$ is a sequence of maps $f_n:G\to\Sym(m_n)$ that satisfy
$$\lim_{n\to\infty}d_{m_n}(f_n(gh),f_n(g)f_n(h))=0$$
for all $g,h\in G$, and 
$$\lim_{n\to\infty}d_{m_n}(f_n(g),id)=1$$
for every $g\neq 1$. In other words, a sofic approximation of $G$ is a separating almost-homomorphism of $G$.

The group $G$ is said to be \emph{sofic} if it admits a sofic approximation. Equivalently, $G$ is sofic if and only if there exist integers $m_n$, a nonprincipal ultrafilter
$\mathcal U$ on $\mathbb N$, and a homomorphism
$\Theta:G\to\mathfrak{S}_\mathcal{U}$
such that
$d_{\mathcal U}\bigl(\Theta(g),1\bigr)=1$
for every $g\neq1$.

% Let $(f_n)_n$ be a sofic approximation of $G$. Choose a nonprincipal ultrafilter $\mathcal{U}$ and form the metric ultraproduct
% $$\mathfrak{S}_{\mathcal{U}}=\prod_{\mathcal{U}} (\Sym(m_n),d_{m_n}),$$
% see \cite{Gol22} for details. Define $f:G\to\mathfrak{S}_{\mathcal{U}}$ by
% $$f(g)=[f_n(g)]_{\mathcal{U}}.$$
% The fact that $(f_n)_n$ is an almost-homomorphism implies that $f$ is a homomorphism of groups, and the separation property of $(f_n)_n$ yields $d_{\mathcal{U}}(f(g),1)=1$ for all $g\in G\setminus\{1\}$.
%%%%%%%%%%%%%%%%%%%%%%%%%%%%%%%

%%%%%%%%%%%%%%%%%%%%%%%%%%%%%%%%%%%%%%%%%%%%%%%%%%%%%%%%%%%%%%%%%%%%%%%

\section{Almost-polynomials}
\label{s:almostpolynomials}

\noindent

\begin{definition}[Almost-polynomial]
    A sequence of 
    asymptotically unital 
    maps $f_n:G\to\Sym(m_n)$ is \emph{almost-polynomial of degree $\le s$} if
    $$\lim_{n\to\infty} d_{m_n}((\mypartial_{g_1,\ldots ,g_{s+1}}f_n)(1),1)=0$$
    for all $g_1,\ldots ,g_{s+1}\in G$.
\end{definition}

Almost-polynomials are closely related to uniform $\varepsilon$-polynomials, \cite{JT24}. Let $(H,d)$ be a group with bi-invariant metric. A unital map $\phi:G\to H$ is a \emph{uniform $\varepsilon$-polynomial of degree $\le s$} if
    $$d(\mypartial_{g_1,...,g_{d+1}}\phi(1),1)\le\varepsilon$$
    for all $g_1,\ldots ,g_{s+1}\in G$.  Note, for example, that if $G$ is a finite group and $(f_n)_n$ is a unital almost-polynomial of $G$, then, given $\varepsilon>0$, all but possibly finitely many $f_n$ are uniform $\varepsilon$-polynomials $G\to\Sym(m_n)$.

\begin{example}
    \label{ex:deg1}
    Almost-polynomials of degree $1$ are precisely almost-homomorphisms. Assume first the sequence of maps $f_n:G\to\Sym(m_n)$ is an almost-homomorphism. Then $(f_n)_n$ is clearly asymptotically unital and
    \begin{align*}
        d_{m_n}((\mypartial_{h,g}f_n)(1),1) &=
        d_{m_n}(f_n(gh)f_n(h)^{-1}f_n(1)f_n(g)^{-1},1)\\
        %&= d_{m_n}(f_n(gh)f_n(h)^{-1}f_n(g)^{-1}\cdot f_n(g)f_n(1)f_n(g)^{-1},1)\\
        &= d_{m_n}(f_n(gh)f_n(h)^{-1}f_n(g)^{-1},f_n(g)f_n(1)^{-1}f_n(g)^{-1})\\
        &\le d_{m_n}(f_n(gh)f_n(h)^{-1}f_n(g)^{-1},1)\\
        &\quad +
        d_{m_n}(1, f_n(g)f_n(1)^{-1}f_n(g)^{-1})\\
        &= d_{m_n}(f_n(gh),f_n(g)f_n(h))+d_{m_n}(f_n(1),1),
    \end{align*}
    and the latter converges to $0$ for all $g,h\in G$. 

    Conversely, let the maps $f_n:G\to\Sym(m_n)$ form an almost-polynomial of $G$ of degree $\le 1$. Then
    \begin{align*}
        d_{m_n}(f_n(gh),f_n(g)f_n(h)) 
        %&=d_{m_n}(f_n(gh)f_n(h)^{-1}f_n(1)f_n(g)^{-1},f_n(g)f_n(1)f_n(g)^{-1})\\
        &= d_{m_n}((\mypartial_{h,g}f_n)(1),f_n(g)f_n(1)f_n(g)^{-1})\\
        &\le d_{m_n}((\mypartial_{h,g}f_n)(1),1)\\
        &\quad +d_{m_n}(1,f_n(g)f_n(1)f_n(g)^{-1})\\
        &=d_{m_n}((\mypartial_{h,g}f_n)(1),1)+d_{m_n}(1,f_n(1)),
    \end{align*}
    and the latter converges to $0$ for all $g,h\in G$.
\end{example}

\begin{proposition}
    \label{prop:degreeup}
    Let a sequence of maps $f_n:G\to\Sym(m_n)$ be an almost-polynomial of degree $\le s$. Then it is also almost-polynomial of degree $\le s+1$.
\end{proposition}

\begin{proof}
    Let $g_1,\ldots ,g_{s+1},g_{s+2}\in G$. Let $\phi_n=\mypartial_{g_2,\ldots ,g_{s+2}}f_n$. The map $\phi_n$ is asymptotically unital. We have 
    \begin{align*}
        d_{m_n}((\mypartial_{g_1,\ldots ,g_{s+2}}f_n)(1),1)
        &= 
        d_{m_n}((\mypartial_{g_1}\phi_n)(1),1)\\
        &= d_{m_n}(\phi_n(g_1)\phi_n(1)^{-1},1)\\
        &\le d_{m_n}(\phi_n(g_1),1)+d_{m_n}(\phi_n(1),1).
    \end{align*}
    The second summand of the latter sum tends to $0$. Also note that (Subsection \ref{ss:poly_prelim})
    \begin{align*}
        \phi_n(g_1) &= (\mypartial_{g_2,\ldots ,g_{s+2}}f_n)(g_1)\\
        &= (\mypartial_{g_2g_1,g_3,\ldots ,g_{s+2}}f_n)(1)
        (\mypartial_{g_1,g_3,\ldots ,g_{s+2}}f_n)(1)^{-1}
    \end{align*}
    and so the first summand goes to $0$ as well.
\end{proof}

\begin{lemma}
    \label{lem:lowerdeg}
    Let $f:G\to H$ be a unital polynomial of degree $\le s$, $s\ge 2$. Then there exists a normal subgroup $K$ of $G$ such that $G/K$ is abelian and $f\big|_K$ is a unital polynomial of degree $\le s-1$.
\end{lemma}

\begin{proof}
    Fix $\mathbf{a}=(a_2,\ldots ,a_s)\in G^{s-1}$. Denote $q_{\mathbf{a}}=\mypartial_{a_2,\ldots ,a_s}f$. As $f$ is polynomial of degree $\le s$, the map
    $\lambda_{\mathbf{a}}(x)=q_{\mathbf{a}}(x)q_{\mathbf{a}}(1)^{-1}$ is a homomorphism \cite[Lemma 2.1]{JT24}. We claim that the range of $\lambda_{\mathbf{a}}$ is abelian. Denote $\mu=\mypartial_{a_3,\ldots ,a_s}f$, and note that $q_{\mathbf{a}}=\mypartial_{a_2}\mu$. Setting $\mu^*(x)=\mu(x)\mu(1)^{-1}$, we have that $\mu^*$ is a unital quadratic map, and $q_{\mathbf{a}}=\mypartial_{a_2}\mu=\mypartial_{a_2}\mu^*$. Fix $k\in G$ and let $\beta_k(x)=\mu^*(k)^{-1}\mypartial_k\mu^*(x)$. By \cite[Proposition 2.4]{JT24}, the group $\beta_k(G)$ is abelian for every $k\in G$. Observe that
    $$\beta_{a_2}(x)=\mu^*(a_2)^{-1}q_{\mathbf{a}}(x)$$
    and 
    $$q_{\mathbf{a}}(1)=\mypartial_{a_2}\mu^*(1)=\mu^*(a_2)\mu^*(1)^{-1}=\mu^*(a_2).$$ 
    These two yield
    $$\lambda_{\mathbf{a}}(x)=\mu^*(a_2)\beta_{a_2}(x)\mu^*(a_2)^{-1},$$
    hence $\lambda_{\mathbf{a}}(G)=\mu^*(a_2)\beta_{a_2}(G)\mu^*(a_2)^{-1}$
    is abelian.

    Set
    $$K=\bigcap_{\mathbf{a}\in G^{s-1}}\ker\lambda_{\mathbf{a}}.$$
    The diagonal map
    $$G\to \prod_{\mathbf{a}\in G^{s-1}}\lambda_{\mathbf{a}}(G)$$
    has kernel $K$ and abelian target. This proves that $G/K$ is abelian.

    Pick $k_1,\ldots ,k_s\in K$. Then
    \begin{align*}
        \mypartial_{k_1,\ldots ,k_s}f(1) &= \mypartial_{k_1}q_{(k_2,\ldots ,k_s)}(1)\\
        &= q_{(k_2,\ldots ,k_s)}(k_1)q_{(k_2,\ldots ,k_s)}(1)^{-1}\\
        &= \lambda_{(k_2,\ldots ,k_s)}(k_1)\\
        &= 1,
    \end{align*}
    where the last step follows from the fact that $k_1\in K\subset\ker \lambda_{(k_2,\ldots ,k_s)}$.
\end{proof}

\begin{corollary}
    \label{cor:degreelowalmost}
    Let $(f_n)_n$ be an almost-polynomial of $G$ of degree $\le s$, where $s\ge 2$. Then $(f_n\big|_{G'})_n$ is an almost-polynomial of $G'$ of degree $\le s-1$.
\end{corollary}

\begin{proof}
    Let $(f_n)_n$ be an almost-polynomial of $G$ of degree $\le s$, $s\ge 2$. Pick an arbitrary nonprincipal ultrafilter $\mathcal{U}$, and form the metric ultraproduct
    $$\mathfrak{S}_{\mathcal{U}}=\prod_{\mathcal{U}} (\Sym(n),d_{m_n}),$$
    Define $f:G\to\mathfrak{S}_{\mathcal{U}}$ by
    $$f(g)=[f_n(g)]_{\mathcal{U}}.$$ By definition of $\mathfrak{S}_{\mathcal{U}}$, the asymptotic identities 
    $$d_{m_n}((\mypartial_{g_1,\ldots ,g_{s+1}}f_n)(1),1)\to 0$$
     yield
    $$\mypartial_{g_1,\ldots ,g_{s+1}}f(1)=1$$
    for all $g_1,\ldots ,g_{s+1}\in G$. Therefore $f$ is a genuine unital polynomial of degree $\le s$.
     By Lemma \ref{lem:lowerdeg}, the restriction of $f$ to $G'$ is a polynomial of degree $\le s-1$, that is,
    $$\lim_{n\to\mathcal{U}}d_{m_n}
    ((\partial_{g_1,\dots,g_s}f_n)(1),1)=0$$
    for all $g_1,\ldots ,g_s\in G$.
    As this is true for every nonprincipal ultrafilter $\mathcal{U}$, we conclude that
    $$\lim_{n\to\infty}d_{m_n}
    ((\partial_{g_1,\dots,g_s}f_n)(1),1)=0,$$
    and the result follows.
\end{proof}

The next consequence can be compared with \cite[Theorem 5.5]{JT24}:
\begin{corollary}
    \label{cor:perfect}
    Let $G$ be a perfect group and let $(f_n)_n$ be an almost-polynomial of $G$. Then $(f_n)_n$ is an almost-homomorphism of $G$.
\end{corollary}

%%%%%%%%%%%%%%%%%%%%%%%%%%%%%%%%%%%%%%%%%%%%%%%%%%%%%%%%%%%%%%%%%%%%%%%

\section{Polynomial permutation stability}
\label{s:polystability}

\noindent

\begin{definition}[Close to polynomial]
    A sequence of maps $f_n:G\to\Sym(m_n)$ is \emph{close to a polynomial of degree $\le s$} if there exists a sequence of unital polynomials $\rho_n:G\to \Sym(m_n)$ of degree $\le s$ such that
    $$\lim_{n\to \infty} d_{m_n}(\rho_n(g),f_n(g))=0$$
for all $g\in G$.
\end{definition}


\begin{definition}[Permutation stable of degree $s$]
    A group $G$ is \emph{permutation stable of degree $\le s$} (\emph{$P_s$-stable} for short) if every almost-polynomial of degree $\le s$ of $G$ is close to a polynomial of degree $\le s$.
\end{definition}

% \begin{remark}
%     \label{rem:almostpolyunital}
%     Suppose $f_n:G\to\Sym(m_n)$ is an almost-polynomial of degree $\le s$. Define $f_n^*(x)=f_n(x)f_n(1)^{-1}$. Then the maps $f_n^*$ are unital. As $\mypartial_gf_n=\mypartial_gf_n^*$, the maps $f_n^*$ also form an almost-polynomial of degree $\le s$. Suppose that $f_n^*$ is close to a polynomial $\rho_n^*:G\to\Sym(m_n)$ of degree $\le s$. Then $f_n$ is close to the polynomial $\rho_n$ of degree $\le s$ defined by $\rho_n(g)=\rho_n^*(g)f_n(1)$. This shows that, regarding the questions of $P_s$-stability, we may assume without loss of generality that the almost-polynomials in question are unital. 
% \end{remark}



\begin{proof}[Proof of Theorem \ref{thm:Pscharacterization}]
    We assume the notations set in Subsection \ref{ss:universal_prelim}.

    Let $G$ be a $P_s$-stable group. Let $f_n:\Pold_s(G)\to\Sym(m_n)$ be an almost-homomorphism of $\Pold_s(G)$. Define $\phi_n:G\to\Sym(m_n)$ by the rule $\phi_n=f_n\circ \iota_s$. As $\iota_s$ is unital and $(f_n)_n$ is an almost-homomorphism, it follows that $(\phi_n)_n$ is asymptotically unital. Fix $\mathbf{g}=(g_1,\ldots ,g_{s+1})\in G^{s+1}$. Then
    \begin{align*}
    d_{m_n}((\mypartial_{\mathbf{g}}\phi_n)(1),1) &\le
    d_{m_n}((\mypartial_{\mathbf{g}}(f_n\circ \iota_s))(1),
    f_n((\mypartial_{\mathbf{g}}\iota_s)(1)))\\
    &\quad + d_{m_n}(f_n((\mypartial_{\mathbf{g}}\iota_s)(1)),1). 
    \end{align*}
    As $\iota_s$ is a polynomial of degree $\le s$ and $(f_n)_n$ is an almost-homomorphism, we have that
    the latter sum converges to 0, therefore $(\phi_n)_n$ is an almost-polynomial of degree $\le s$ of $G$. By the $P_s$-stability of $G$ there exist unital polynomials $\rho_n:G\to\Sym(m_n)$ of degree $\le s$ with
    $$\lim_{n\to\infty}d_{m_n}(\phi_n(g),\rho_n(g))=0$$
    for every $g\in G$. By Proposition \ref{prop:Pold}, there exist (uniquely determined) homomorphisms $\psi_n:\Pold_s(G)\to\Sym(m_n)$ with $\rho_n=\psi_n\circ\iota_s$. Take $a\in\Pold_s(G)$ and write it as
$$a=\prod_{i=1}^k\iota_s(g_i)^{\varepsilon_i},$$
where $g_i\in G$. and $\varepsilon_i=\pm 1$. Then
$$\psi_n(a)=\prod_{i=1}^k\psi_n(\iota_s(g_i))^{\varepsilon_i}=\prod_{i=1}^{k}\rho_n(g_i)^{\varepsilon_i}.$$
On the other hand,
%$$f_n(a)=f_n\left (\prod_{i=1}^{k}\iota_s(g_i)^{\varepsilon_i}\right )$$
%is close \comment{(DETAILS!!)} to 
%$$\prod_{i=1}^kf_n(\iota_s(g_i))^{\varepsilon_i}=\prod_{i=1}^k\phi_n(g_i)^{\varepsilon_i},$$
$$d_{m_n}\left (f_n(a),\prod_{i=1}^k\phi_n(g_i)^{\varepsilon_i}\right )=d_{m_n}\left (f_n\left (\prod_{i=1}^{k}\iota_s(g_i)^{\varepsilon_i}\right ), \prod_{i=1}^kf_n(\iota_s(g_i))^{\varepsilon_i}\right )$$
converges to 0, hence
\begin{align*}
    d_{m_n}(\psi_n(a),f_n(a)) &\le d_{m_n}\left ( \prod_{i=1}^{k}\rho_n(g_i)^{\varepsilon_i}, \prod_{i=1}^k\phi_n(g_i)^{\varepsilon_i}\right )+
    d_{m_n}\left (f_n(a),\prod_{i=1}^k\phi_n(g_i)^{\varepsilon_i}\right )\\
    &\le \sum_{i=1}^k d_{m_n}(\rho_n(g_i),\phi_n(g_i))+d_{m_n}\left (f_n(a),\prod_{i=1}^k\phi_n(g_i)^{\varepsilon_i}\right )
\end{align*}
converges to 0.

Conversely, assume that $\Pold_s(G)$ is $P$-stable. Let the sequence $f_n:G\to\Sym(m_n)$ be an almost-polynomial of $G$ of degree $\le s$. Let $\widetilde{f}_n:F_G\to\Sym(m_n)$ be the unique homomorphism mapping $x_g$ to $f_n(g)$ for all $g\in G$. By the assumption on $f_n$ we have that
$$\lim_{n\to\infty} d_{m_n}(\widetilde{f}_n((\mypartial_{g_1,\ldots ,g_{s+1}}\iota)(1)),1)=0$$
for all $g_1,\ldots ,g_{s+1}\in G$. By Lemma \ref{lem:Pstabrel}, there exist homomorphisms $\rho_n:\Pold_s(G)\to\Sym(m_n)$ such that
$$\lim_{n\to\infty}d_{m_n}(\widetilde{f}_n(w),\rho_n(wR_s(G)))=0$$
for every $w\in F$. The maps $\phi_n:G\to\Sym(m_n)$, $\phi_n=\rho_n\circ\iota_s$, are unital polynomials of degree $\le s$, and
$$d_{m_n}(f_n(g),\phi_n(g))=d_{m_n}(\widetilde{f}_n(x_g),\rho_n(\iota_s(g)))=
d_{m_n}(\widetilde{f}_n(x_g),\rho_n(x_gR_s(G)))$$
converges to zero for every $g\in G$. This proves the claim.
\end{proof}

\begin{corollary}
    \label{cor:finiteP2}
    Finite groups are $P_2$-stable.
\end{corollary}

\begin{proof}
    If $G$ is finite, then $\Pold_2(G)$ is finite by \cite[Corollary 4.4]{JT24}, hence $\Pold_2(G)$ is $P$-stable by \cite{GR09}.
\end{proof}

We record another consequence of a stability criterion of Arzhantseva and
P\u{a}unescu \cite{AP24}. Recall that a subgroup $H$ of a group $\Gamma$ is
\emph{almost normal} if it has only finitely many conjugates, or
equivalently,
$[\Gamma:N_\Gamma(H)]<\infty$.
We also say that $H$ is \emph{profinitely closed} if
$H=\bigcap_{i=1}^{\infty} H_i$
for some finite-index subgroups $H_i\leq \Gamma$.
The following is a special case of the results of
Arzhantseva--P\u{a}unescu on action traces and constraint stability.

\begin{proposition}[\cite{AP24}]
\label{prop:AP-criterion}
Let $\Gamma$ be a countable amenable group.  Suppose that
$\Sub(\Gamma)$ is countable and that every almost normal subgroup of
$\Gamma$ is profinitely closed.  Then $\Gamma$ is $P$-stable.
\end{proposition}

\begin{proof}
Arzhantseva and P\u{a}unescu prove that, under the assumptions that
$\Sub(\Gamma)$ is countable and every almost normal subgroup is
profinitely closed, the relevant action traces are residually finite,
see \cite[Proposition 5.17]{AP24}.  Combining this with amenability
gives constraint stability in
\cite[Corollary~ 5.19]{AP24}.  Taking the constraint subgroup to be
trivial yields ordinary $P$-stability.
\end{proof}

We obtain the following polynomial consequence.

\begin{corollary}
\label{cor:AP-polynomial}
Let $G$ be a countable group and let $s\geq1$.  Suppose that
$P=\Pold_s(G)$
is amenable, that $\Sub(P)$ is countable, and that every almost normal subgroup of $P$ is profinitely closed.  Then $G$ is $P_s$-stable.
\end{corollary}

%%%%%%%%%%%%%%%%%%%%%%%%%%%%%%%%%%%%%%%%%%%%%%%%%%%%%%%%%%%%%%%%%%%%%%%

\section{Polynomially sofic approximations}
\label{s:polysoficapprox}

\noindent

\begin{definition}[Polynomially sofic approximation]
    Let $G$ be a countable group.
    The maps $f_n:G\to\Sym(m_n)$ form a \emph{polynomial sofic approximation of $G$ of degree $\le s$} if $(f_n)_n$ is an almost-polynomial of degree $\le s$ of $G$ that is separating, i.e., 
    $$\lim_{n\to\infty} d_{m_n}(f_n(g),id)=1$$
    for all $g\in G\setminus\{1\}$.
\end{definition}

\begin{remark}
    \label{rem:unitalsof}
    Recall that almost-homomorphisms are asymptotically unital by definition. Without assuming this, the notion of polynomial sofic approximations becomes degenerate. Namely, take any sequence of permutations $\sigma_n\in\Sym(m_n)$ with
    $$\lim_{n\to \infty }d_{m_n}(\sigma_n,1) =1$$
    and consider the constant maps
    $f_n(g)=\sigma_n$, $g\in G$. These are genuine polynomials of any degree, and
    $$\lim_{n\to\infty}d_{m_n}(f_n(g),1)=d_{m_n}(\sigma_n,1)=1.$$
    Thus every group would be admitting a polynomial sofic approximation of any degree under that definition.
\end{remark}

\begin{proof}[Proof of Theorem \ref{thm:soficcharacterization}]
    If $G$ is sofic, it admits a polynomial sofic approximation of degree $\le 1$, which is a polynomial sofic approximation of degree $\le s$ for every $s\ge 1$.

    The converse is proved by induction on $s$. Let $s\ge 2$ and let $G$ admit a polynomial sofic approximation $f_n:G\to\Sym(m_n)$ of degree $\le s$. Pick an arbitrary nonprincipal ultrafilter $\mathcal{U}$, and form the metric ultraproduct
    $\mathfrak{S}_{\mathcal{U}}$.
    Then $(f_n)_n$ induces a unital polynomial $f:G\to\mathfrak{S}_{\mathcal{U}}$ of degree $\le s$.
    Lemma \ref{lem:lowerdeg} implies that the map $f\mid_{G'}$ is a polynomial of degree $\le s-1$, hence it induces a polynomial sofic approximation of $G'$ of degree $s-1$. By induction, $G'$ is sofic, therefore $G$ is sofic.
\end{proof}

\begin{corollary}
    \label{cor:soficG}
    Let $G$ be a countable group. If $\Pold_s(G)$ is sofic, then $G$ is sofic.
\end{corollary}

\begin{proof}
    Let $\mathcal{U}$ be a nonprincipal ultrafilter admitting a homomorphism $f:\Pold_s(G)\to\mathfrak{S}_{\mathcal{U}}$ with $d_{\mathcal{U}}(f(p),1)=1$ for every $p\in\Pold_s(G)\setminus\{1\}$. Let $\phi=f\circ\iota_s:G\to\mathfrak{S}_{\mathcal{U}}$. Then $\phi$ is a unital polynomial of degree $\le s$. As $\iota_s$ is injective, we have that $d_{\mathcal{U}}(\phi(g),1)=1$ for every $g\in\setminus\{1\}$. This implies that $\phi$ induces a polynomial sofic approximation of $G$ of degree $\le s$. By Theorem \ref{thm:soficcharacterization}, $G$ is sofic.
\end{proof}

\begin{proposition}
    \label{prop:soficPol2}
    Let $G$ be a sofic group. Then $\Pold_2(G)$ is sofic.
\end{proposition}

\begin{proof}
    It follows from \cite[Theorem 1.3]{JT24} that $\Pold_2(G)$ is isomorphic to the twisted cross product $A\rtimes_\psi G$, where $A=\omega(G)\otimes_\mathbb{Z} G^{\rm ab}$. Here $\omega(G)$ is the augmentation ideal of $\mathbb{Z}G$. Note that the embedding of $\omega(G)$ into $\mathbb{Z}G$ is split, hence tensoring with $G^{\rm ab}$ keeps it injective.
    The 2-cocycle $\psi:G\times G\to A$ is given by $\psi(g,h)=(g-1)\otimes \bar{h}$, where $\bar{h}=hG'$.

    Consider $M=\mathbb{Z}G\otimes_\mathbb{Z} G^{\rm ab}$. As a $G$-module,
    $$M\cong \bigoplus_{g\in G} G^{\rm ab},$$
    with $G$ acting by left translation. Thus $M\rtimes G\cong G^{\rm ab}\wr G$. Define $\Phi:A\rtimes_\psi G\to M\rtimes G$ by
    $$\Phi(\xi,g)=(\xi+1\otimes \bar{g},g)$$
    for all $\xi\in A, g\in G$. As
    \begin{align*}
        \Phi((\xi_1,g_1),(\xi_2,g_2)) &= 
        \Phi(\xi_1+g_1\xi_2+\psi(g_1,g_2),g_1g_2)\\
        &= (\xi_1+g_1\xi_2+\psi(g_1,g_2)+1\otimes (\bar{g}_1+\bar{g}_2),g_1g_2)\\
        &= (\xi_1+g_1\xi_2+g_1\otimes \bar{g}_2+1\otimes \bar{g}_1,g_1g_2)\\
        &=\Phi(\xi_1,g_1)\Phi(\xi_2,g_2),
    \end{align*}
    it follows that $\Phi$ is an injective homomorphism, hence $\Pold_2(G)$ embeds in $G^{\rm ab}\wr G$. By \cite[Theorem 1.1]{HS18}, $\Pold_2(G)$ is sofic.
\end{proof}

\begin{example}[Polynomial sofic approximation of degree 2 for $\mathbb{Z}$]
    \label{ex:polysofZ}
    Let $p_n$ be a sequence of strictly increasing odd primes. Let $X_n=C_{p_n}\times C_{p_n}$ and $m_n=|X_n|=p_n^2$. For $x,y\in X_n$ let $a_n(x,y)=(x+1,y)$ and $b_n(x,y)=(x,y+1)$. The map $q_n:\mathbb{Z}\to\Sym(X_n)$ given by 
    $$q_n(k)=a_n^kb_n^{{k\choose 2}}$$
    is a unital quadratic map. Let $\tau_n\in\Sym(X_n)$ be the transposition interchanging $(0,0)$ and $(0,1)$. Define $f_n:\mathbb{Z}\to\Sym(X_n)$ by the rule
    $$f_n(k)=\left\{ \begin{array}{ccc}
        q_n(k) & : & k\neq 2\\
        \tau_n\cdot q_n(2) & : & k=2
    \end{array}
    \right . .$$
    For a fixed $k\in\mathbb{Z}$ we have that
    $$d_{m_n}(f_n(k),q_n(k))\le \frac{2}{p_n^2}\rightarrow 0.$$
    As the Hamming distance is bi-invariant, we therefore get
    $$d_{m_n}((\mypartial_{k_1,k_2,k_3}f_n)(1), 1)=
    d_{m_n}((\mypartial_{k_1,k_2,k_3}f_n)(1),\mypartial_{k_1,k_2,k_3}q_n)(1))\rightarrow 0$$
    for all $k_1,k_2,k_3\in\mathbb{Z}$. This shows that $(f_n)_n$ is an almost-polynomial of $\mathbb{Z}$ of degree exactly 2. Note that $f_n$ are not quadratic maps as $(\mypartial_{1,1,1}f_n)(0)\neq 1$. Fix a non-zero integer $k$. For sufficiently large $n$ we have that $p_n>|k|$, so $k\not\equiv 0\mod p_n$, therefore $q_n(k)$ is a nontrivial translation on $X_n$. Therefore we have that $d_{m_n}(q_n(k),1)=1$ for all $n$ large enough. In this case,
    $$d_{m_n}(f_n(k),1)\ge d_{m_n}(q_n(k),1)-d_{m_n}(f_n(k),q_n(k))\ge 1-\frac{2}{p_n^2}\rightarrow 1.$$
    This shows that $(f_n)$ is a sofic approximation of degree $2$ for $\mathbb{Z}$.
\end{example}

%%%%%%%%%%%%%%%%%%%%%%%%%%%%%%%%%%%%%%%%%%%%%%%%%%%%%%%%%%%%%%%%%%%%%%%

\section{Polynomial weak permutation stability of finite groups}
\label{s:weakpolypermut}

\noindent
We next introduce a weak version of polynomial permutation stability,
following the notion of weak stability of Arzhantseva and
P\u{a}unescu \cite[Section7]{AP15}. We adopt the notations of Subsection \ref{ss:universal_prelim}.

Throughout this section, let $G$ be a finite group.  
Given a map
$f:G\longrightarrow\Sym(m)$,
we denote by
$\widehat f:F_G\longrightarrow\Sym(m)$
the unique homomorphism satisfying
$\widehat f(x_g)=f(g)$ for all $g\in G$. 
Thus $f$ is an exact unital polynomial map of degree at most $s\ge 1$ if and only if
$R_s(G)\leq\ker\widehat f$.

Let $\ell$ denote the word length on $F_G$ with respect to the
generating set $\{x_g:g\in G\}$.
Let $\delta>0$.  A map
$f:G\longrightarrow\Sym(m)$
is called a \emph{$\delta$-strong polynomial approximation of degree
at most $s$} if, for every $w\in F_G$ with
$\ell(w)<1/\delta$,
one has
$$w\in R_s(G)
\quad\Longrightarrow\quad
d_m\bigl(\widehat f(w),1\bigr)<\delta,$$
and
$$w\notin R_s(G)
\quad\Longrightarrow\quad
d_m\bigl(\widehat f(w),1\bigr)>1-\delta.$$
This is precisely the notion of a $\delta$-strong solution introduced
in \cite[Definition~7.1]{AP15}, applied to the canonical presentation
$\Pold_s(G)=F_G/R_s(G)$.

\begin{definition}
\label{def:weak-Ps-stability}
We say that $G$ is \emph{weakly $P_s$-stable} if, for every
$\varepsilon>0$, there exists $\delta>0$ such that every
$\delta$-strong polynomial approximation$
f:G\longrightarrow\Sym(m)$
of degree at most $s$ is $\varepsilon$-close to an exact unital
polynomial map
$q:G\longrightarrow\Sym(m)$
of degree $\le s$,
that is,
$$\max_{g\in G}d_m\bigl(f(g),q(g)\bigr)<\varepsilon.$$
\end{definition}

The terminology is justified by the following immediate consequence
of the universal property.

\begin{proposition}
\label{prop:weak-Ps-universal}
Let $G$ be finite and $s\geq1$.  Then
$G$ weakly $P_s$-stable if and only if
$\Pold_s(G)$ is weakly stable in permutations.
\end{proposition}

% AI!!!!!
% \begin{proof}
% The canonical presentation of $\Pold_s(G)$ is
% \[
% \Pold_s(G)=F_G/N_s.
% \]
% By Definition~\ref{def:strong-polynomial}, a
% $\delta$-strong polynomial approximation
% \[
% f:G\to\Sym(m)
% \]
% is exactly a $\delta$-strong solution, in the sense of
% \cite[Definition~7.1]{AP15}, of this presentation, under the
% identification
% \[
% x_g\longmapsto f(g).
% \]

% Moreover, an exact solution of the presentation is precisely a
% homomorphism
% \[
% \rho:\Pold_s(G)\longrightarrow\Sym(m).
% \]
% By the universal property of $\Pold_s(G)$, such homomorphisms are in
% bijection with exact unital polynomial maps
% \[
% q:G\longrightarrow\Sym(m),
% \qquad
% q=\rho\circ\iota_s,
% \qquad
% \deg(q)\leq s.
% \]
% Finally, closeness of the two solutions on the generators $x_g$ is
% exactly the condition
% \[
% \max_{g\in G}d_m(f(g),q(g))<\varepsilon.
% \]
% Thus Definition~\ref{def:weak-Ps-stability} is precisely weak
% stability of the canonical presentation of $\Pold_s(G)$.
% \end{proof}
% End AI

For our use it is convenient to record the corresponding sequential
formulation.

\begin{definition}
A sequence of maps
$f_n:G\longrightarrow\Sym(m_n)$
is called a \emph{strong polynomial approximation of degree at most
$s$} if, for every $w\in F_G$,
\[
w\in R_s(G)
\quad\Longrightarrow\quad
d_{m_n}\bigl(\widehat f_n(w),1\bigr)\longrightarrow0,
\]
whereas
\[
w\notin R_s(G)
\quad\Longrightarrow\quad
d_{m_n}\bigl(\widehat f_n(w),1\bigr)\longrightarrow1.
\]
\end{definition}

Thus the first condition says that all relations of $\Pold_s(G)$ are
satisfied asymptotically, while the second says that all non-relations
are asymptotically maximally separated from the identity.

\begin{proposition}
\label{prop:weak-Ps-sequential}
The group $G$ is weakly $P_s$-stable if and only if, for every strong
polynomial approximation
$f_n:G\longrightarrow\Sym(m_n)$
of degree $\le s$, there exist exact unital polynomial maps
$q_n:G\longrightarrow\Sym(m_n)$
of degree $\le s$
such that
$$
\max_{g\in G}
d_{m_n}\bigl(f_n(g),q_n(g)\bigr)
\longrightarrow0.
$$
\end{proposition}

\begin{proof}
Suppose first that $G$ is weakly $P_s$-stable.  If $(f_n)$ is a
strong polynomial approximation, then for every fixed $\delta>0$ the
set
\[
\{w\in F_G:\ell(w)<1/\delta\}
\]
is finite.  Hence, for all sufficiently large $n$, the map $f_n$ is a
$\delta$-strong polynomial approximation.  Applying weak
$P_s$-stability with $\delta$ corresponding to a sequence
$\varepsilon\to0$ gives the required polynomial maps $q_n$.

Conversely, if weak $P_s$-stability fails, there exist
$\varepsilon>0$, a sequence $\delta_n\to0$, and
$\delta_n$-strong polynomial approximations
\[
f_n:G\to\Sym(m_n)
\]
which remain at distance at least $\varepsilon$ from every exact
degree-$s$ polynomial map.  For each fixed $w\in F_G$, eventually
\[
\ell(w)<1/\delta_n,
\]
and therefore the defining strong-solution inequalities show that
$(f_n)$ is a strong polynomial approximation in the sequential
sense.  This contradicts the assumed sequential correction property.
\end{proof}
Arzhantseva and P\u{a}unescu prove that a finitely presented group
$\Gamma$ is weakly stable if and only if every sofic representation of
$\Gamma$ into a metric ultraproduct of symmetric groups is perfect;
see \cite[Definition 4.1 and Theorem 7.2(i)]{AP15}.
Applied to the universal polynomial group, their result gives the
following formulation:

\begin{corollary}
\label{cor:weak-Ps-sofic-representations}
Let $G$ be finite and $s\geq1$.  Then $G$ is weakly $P_s$-stable if
and only if every sofic representation
$$\Theta:
\Pold_s(G)\longrightarrow
\prod_{n\to\mathcal U}
\bigl(\Sym(m_n),d_{m_n}\bigr)$$
is perfect.
\end{corollary}

% Equivalently, a strong polynomial approximation of $G$ is simply a
% sofic approximation of $\Pold_s(G)$ expressed in terms of the
% distinguished generating set $\iota_s(G)$.

% Indeed, if
% \[
% f_n:G\to\Sym(m_n)
% \]
% is strong, choose a section
% \[
% \tau:\Pold_s(G)\to F_G
% \]
% of $\pi_s$ and define
% \[
% \sigma_n(p)=\widehat f_n(\tau(p)).
% \]
% The defining conditions imply that $(\sigma_n)$ is a sofic
% approximation of $\Pold_s(G)$.  Conversely, restricting any sofic
% approximation of $\Pold_s(G)$ to $\iota_s(G)$ gives a strong polynomial
% approximation.

%\subsection{The amenable case}

The preceding formulation becomes particularly simple when the
universal polynomial group is amenable.

\begin{corollary}
\label{cor:weak-Ps-amenable}
Let $G$ be finite and $s\geq1$, and suppose that $\Pold_s(G)$ is
amenable.  Then
$G$ is weakly $P_s$-stable if and only if
$\Pold_s(G)$ is residually finite.
\end{corollary}

\begin{proof}
By Proposition~\ref{prop:weak-Ps-universal}, $G$ is weakly
$P_s$-stable if and only if $\Pold_s(G)$ is weakly stable.
Arzhantseva and P\u{a}unescu prove that an amenable group is weakly
stable if and only if it is residually finite; see
\cite[Theorem~7.2(iii)]{AP15}.
\end{proof}
More generally, without the amenability assumption, one obtains from
\cite[Theorem~7.2(ii)]{AP15} that if
$\Pold_s(G)$ is sofic and $G$ is weakly $P_s$-stable, then
$\Pold_s(G)$ is residually finite.

%%%%%%%%%%%%%%%%%%%%%%%%%%%%%%%%%%%%%%%%%%%%%%%%%%%%%%%%%%%%%%%%%%%%%%%

\section{$\Pold_s(G)$ and largeness}
\label{s:poldfinite}

\noindent
We assume throughout this section that $G$ is finite group. We first record some basic observations. Let $\alpha:G\to Q$ be a homomorphism and  $\iota_s^Q:Q\to\Pold_s(Q)$  and $\iota_s^G:G\to\Pold_s(G)$ the universal polynomials of degree $\le s$ on $Q$ and $G$, respectively. Then the universal property implies that there exists a uniquely determined homomorphism $\alpha_s^*:\Pold_s(G)\to\Pold_s(Q)$ with $\alpha_s^*\circ \iota_s^G=\iota_s\circ \alpha$. This is the functoriality mentioned in \cite{JT24}. Furthermore, if $\alpha$ is surjective, then the image of $\alpha_s^*$ contains $\iota_s^Q(Q)$, which generates $\Pold_s(Q)$, hence $\alpha_s^*$ is also surjective.

\begin{proposition}
    \label{prop:Gabdivisibleby3}
    If $|G^{\rm ab}|$ is divisible by $3$, then $\Pold_s(G)$ is non-amenable for every $s\ge 3$.
\end{proposition}

\begin{proof}
    By the assumption, $G$ has a quotient isomorphic to $C_3$. Note that $\Pold_3(C_3)$ is non-amenable. For every $s\ge 3$, we have an epimorphism
    $$\Pold_s(G)\twoheadrightarrow \Pold_s(C_3)\twoheadrightarrow \Pold_3(C_3),$$
    hence $\Pold_s(G)$ is non-amenable.
\end{proof}

A group $G$ is called \emph{large} if there exists a finite-index
subgroup $H\leq G$ and an epimorphism
$H\twoheadrightarrow F_2$
where $F_2$ denotes the free group of rank two. Note that large groups are clearly non-amenable.

\begin{proposition}
    \label{prop:c2perfect}
    There exist finite groups $G$ with $G'$ perfect, $G^{\rm ab}\cong C_2$, such that $\Pold_s(G)$ are large for all sufficiently large $s$.
\end{proposition}

\begin{proof}
    Let $N$ be a nontrivial finite perfect group of order $q$. Let $\alpha\in\Aut(N)$ be an involution. Form $G=N\rtimes_\alpha \langle t\rangle$, where $\langle t\rangle\cong C_2$. We claim that $\Pold_s(G)$ are nonamenable for sufficiently large $s$.

    Let $f:G\to H$ be a unital polynomial of degree $\le s$, and denote $c=f(t)$. Consider the maps $\mathfrak{a}:N\to H$ and $\mathfrak{b}:N\to H$ given by
    $\mathfrak{a}(n)=f(n)$ and $\mathfrak{b}(n)=f(nt)c^{-1}$. Clearly, $\mathfrak{a}$ and $\mathfrak{b}$ are unital polynomials of degree $\le s$ on $N$. As $N$ is perfect, they are homomorphisms by \cite[]{JT24}.

    Define $\beta_t:G\to H$ by $\beta_t(x)=c^{-1}f(tx)f(x)^{-1}$. This is a normalized left difference of $f$, hence it is a unital polynomial of degree $\le s-1$ on $G$. Consider the maps $\mathfrak{c}:N\to H$ and $\mathfrak{d}:N\to H$ given by $\mathfrak{c}(n)=\beta_t(n)$ and $\mathfrak{d}(n)=\beta_t(nt)\beta_t(t)^{-1}$. These are unital polynomials of degree $\le s-1$ on $N$, hence $\mathfrak{c}$ and $\mathfrak{d}$ are also homomorphisms.
    
    A direct calculation shows that the relation $\alpha(n)t=tn$ yields
    \begin{equation}
        \label{eq:Nt1}
        \mathfrak{c}(n)=c^{-1}\mathfrak{b}(\alpha(n))c\mathfrak{a}(n)^{-1}.
    \end{equation}
    We also note that $\beta_t(t)=c^{-1}f(t^2)f(t)^{-1}=c^{-2}$. This, together with straightforward calculation, shows that
    \begin{equation}
        \label{eq:Nt2}
        \mathfrak{d}(n)=c^{-1}\mathfrak{a}(\alpha(n))c^{-1}\mathfrak{b}(n)^{-1}c^2.
    \end{equation}
    We conclude that
    $$c\mathfrak{d}(\alpha(n))c^{-1}=\mathfrak{a}(\alpha^2(n))c^{-1}\mathfrak{b}(\alpha(n))^{-1}c=\mathfrak{c}(n)^{-1}.$$
    This shows that the map $n\mapsto\mathfrak{c}(n)^{-1}$ is a homomorphism $N\to H$. But $\mathfrak{c}$ is also a homomorphism, therefore we see that $\mathfrak{c}(N)$ needs to be abelian. But this is also a perfect group, hence $\mathfrak{c}(N)=1$, and therefore $\mathfrak{d}(N)=1$.

    The equation \eqref{eq:Nt1} yields $\mathfrak{b}(\alpha(n))=c\mathfrak{a}(n)c^{-1}$, thus $\mathfrak{b}(n)=c\mathfrak{a}(\alpha(n))c^{-1}$. In summary, we have that $f(n)=\mathfrak{a}(n)$ and $f(nt)=c\mathfrak{a}(\alpha(n))$ for all $n\in N$. This implies that 
    \begin{equation}
        \label{eq:Nt3}
        f(gn)=f(g)\mathfrak{a}(n)
    \end{equation}
    for all $g\in G$ and all $n\in N$.

    Given $g\in G$, consider the normalized difference map $\beta_g:G\to H$ given by $\beta_g(x)=f(g)^{-1}f(gx)f(x)^{-1}$. Given $x\in G$ and $n\in N$, we get from Equation \eqref{eq:Nt3} that $\beta_g(xn)=\beta_g(x)$. This shows that $\beta_g$ factors through $G/N\cong C_2$.

    We recall from \cite{JT24,AT26} that $\Pold_s(C_2)\cong C_{2^s}$. In other words, if $p:C_2\to H$ is a unital map sending the generator of $C_2$ to $u\in H$, then $p$ is a polynomial of degree $\le s$ if and only if $u^{2^s}=1$.

    Given $n\in N$, denote
    $$u_n=\beta_n(t)=\mathfrak{a}(n)^{-1}c\mathfrak{a}(\alpha(n))c^{-1}$$
    and
    $$v_n=\beta_{nt}(t)=\mathfrak{a}(\alpha(n))^{-1}c^{-1}\mathfrak{a}(n)c^{-1}.$$
    Note that $f$ is a polynomial of degree $\le s$ if and only $\beta_g$ are polynomials of degree $\le s-1$ for all $g\in G$. By the above, this is precisely when there exist a homomorphism $\mathfrak{a}:N\to H$, and an element $c\in H$ such that
    \begin{align*}
        f(n) &= \mathfrak{a}(n),\\
        f(nt) &= c\mathfrak{a}(\alpha(n)),\\
        u_n^M &= 1,\\
        v_n^M &=1,
    \end{align*}
    for all $n\in N$, where $u_n$ and $v_n$ are as above, and $M=2^{s-1}$. This immediately implies that
    $$\Pold_s(G)\cong \langle N,c\mid (n^{-1}c\alpha(n)c^{-1})^M=(\alpha(n)^{-1}c^{-1}nc^{-1})^M=1,\, n\in N\rangle.$$
    Note that the second set of relators implies $c^{2M}=1$.

   Let $\Gamma_0=N*\langle c\rangle \cong N*\mathbb Z,$
    where at this stage $c$ has infinite order.  The relation
    $c^{2^s}=1$ will be imposed below as one of the
    $2^{s-1}$-power relations. We have an epimorphism $\eta:\Gamma\to G$ sending $n\mapsto n$, $c\mapsto t$. Let $F$ be the kernel of $\eta$. As $\eta$ is injective on $N$, Kurosh's theorem implies that $F$ is free, and
    $|\Gamma:F|=|G|=2q$. By multiplicativity of the group Euler characteristic under passage to
    finite-index subgroups and the free-product formula
    \cite[Chapter IX, Proposition 7.3]{Bro82}, we have
    $$\chi(F)=[\Gamma:F]\chi(\Gamma)=2q\left(\frac1{q}-1\right)=2-2q.$$
    This shows that $F$ is free of rank $2q-1$, hence non-abelian. Note also that the "roots" of all relators of $\Pold_s(G)$ belong to $F$:
    \begin{align*}
        \eta(n^{-1}c\alpha(n)c^{-1}) &=n^{-1}t\alpha(n)t^{-1}=1,\\
        \eta(\alpha(n)^{-1}c^{-1}nc^{-1}) &=\alpha(n)^{-1}tnt=1.
    \end{align*}
    Choose a transversal $\mathcal{T}$ for $F$ in $\Gamma$. As $F$ is a normal subgroup of $\Gamma$, we conclude that the kernel of the epimorphism $\Pold_s(G)\twoheadrightarrow G$ has the form
    $$F_s=F/\llangle (\gamma r\gamma^{-1})^M\mid r\in\mathcal{R},\gamma\in\mathcal{T}\rrangle,$$
    where 
    $$\mathcal{R}=\{ u_n\mid n\in N, n\neq 1\} \cup\{v_n\mid n\in N\}$$
    is a fixed finite set independent of $s$, and the same holds for $\mathcal{T}$. By a result of Lackenby \cite[Theorem 1.2]{Lac07}, $F_s$ is large for sufficiently large $s$. As $F_s$ has finite index in $\Pold_s(G)$, it follows that the latter group is large for sufficiently large $s$. 
    \end{proof}

    The above result gives no particular lower bound for $s$. Keeping the notations of the proof of Proposition \ref{prop:c2perfect}, we have a presentation of $F_s$ with $r=2q-1$ generators and $\le 2q(2q-1)=2qr$ relators, each of which is a  $2^{s-1}$-power of an element in $F$. Hence the $2$-deficiency of $F_s$ is bounded from below as follows:
    $$\deff_2(F_s)\ge r-\frac{2qr}{2^{s-1}}.$$
    It is shown in \cite{BT11} that if the $p$-deficiency of a group is bigger than 1, then the group is large. Thus the largeness of the group $F_s$, and therefore of $\Pold_s(G)$ is ensured by the condition
    $$2^{s-1}>\frac{2qr}{r-1}.$$
    For example, consider the group $\Sym(n)$, where $n\ge 5$. Then we can 
    take $N=\Alt(n)$, and $\Pold_s(\Sym(n))$ is large if
    $$s>1+\log_2\frac{n!(n!-1)}{n!-2}.$$
%%%%%%%%%%%%%%%%%%%%%%%%%%%%%%%%%%%%%%%%%%%%%%%%%%%%%%%%%%%%%%%%%%%%%%%

\section*{Use of AI tools}

\noindent
ChatGPT (GPT-5.6, OpenAI) was used to
assist in drafting parts of this manuscript. All content was reviewed and substantially
revised by the author, who is responsible for the final text.

% VER2
% During the preparation of this manuscript, the author used ChatGPT (GPT-5.6, OpenAI) on and around 14 August 2026, accessed through \url{https://chatgpt.com}. ChatGPT was used solely to assist with the manuscript’s structure, organization, formatting, and presentation. It was not used to conduct the research, analyse data, interpret findings, translate sources, generate images, or develop the manuscript’s substantive arguments or conclusions.
% All suggestions were critically reviewed and revised as necessary by the author. The author independently verified the manuscript’s content, citations, quotations, and references and accepts full responsibility for the accuracy, originality, and integrity of the final work. No confidential, proprietary, or personally identifiable information was submitted to the tool.

%VER3
% The authors used OpenAI tools to assist with basic language editing of portions of this manuscript. The mathematical content, claims, proofs, and conclusions are the authors’ own. The resulting work, in its totality, is an accurate representation of the authors’ underlying work, and the authors take full responsibility for the veracity and correctness of all material.

%%%%%%%%%%%%%%%%%%%%%%%%%%%%%%%%%%%%%%%%%%%%%%%%%%%%%%%%%%%%%%%%%%%%%%%

\begin{thebibliography}{9}

    \bibitem{AT26} 
    V. Alekseev, and A. Thom, 
    Cubic maps from the group of order 3, 
    arXiv:2603.08452, 2026.

    \bibitem{AP15}
    G. Arzhantseva and L. P\u{a}unescu,
    Almost commuting permutations are near commuting permutations,
    J. Funct. Anal. \textbf{269} (2015), no. 3, 745--757.

    \bibitem{AP24}
    G. Arzhantseva and L. P\u{a}unescu,
    Constraint stability in permutations and action traces, Isr. J. Math, in press;
    arXiv:2304.07656v2, 2024.

    \bibitem{BLT19}
    O. Becker, A. Lubotzky, and A. Thom, 
    Stability and invariant random subgroups, 
    Duke Math. J. {\bf 168} (2019), 2207--2234.

    % \bibitem{Bra26}
    % H. Bradford,
    % Local permutation stability,
    % Groups Geom. Dyn. {\bf 20} (2026), 627–-661.

    \bibitem{Bro82}
    K. S. Brown,
    Cohomology of Groups,
    Graduate Texts in Mathematics, vol.~87,
    Springer-Verlag, New York, 1982.

    \bibitem{BT11}
    J. O. Button, and A. Thillaisundaram,
    Applications of $p$-deficiency and $p$-largeness,
    Internat. J. Algebra Comput. {\bf 21} (2011),  547--574.

    \bibitem{GR09}
    L. Glebsky, and L. M. Rivera, 
    Almost solutions of equations in permutations, 
    Taiwanese J. Math. {\bf 13} (2009), 493-500.

    \bibitem{Gol22}
    I. Goldbring, 
    Ultrafilters throughout mathematics, 
    Graduate Studies in Mathematics, Providence, Rhode Island, 2022.

    \bibitem{GT12}
    B. Green, and T. Tao, 
    The quantitative behaviour of polynomial orbits on nilmanifolds, 
    Ann. Math {\bf 175} (2012), 465--540.

    \bibitem{GTZ12}
    B. Green, T. Tao, and T. Ziegler, 
    An inverse theorem for the Gowers $U^{s+1}[N]$-norm,
    Ann. Math. {\bf 176} (2012), 1231–-1372.

    \bibitem{HS18}
    B. Hayes, and A. Sale,
    Metric approximations of wreath products,
    Ann. Inst. Fourier (Grenoble) {\bf 68} (2018), 423--455.

    %\bibitem{Ish26}
    %H. Ishikura, Free metabelian groups are permutation stable,
    %Groups Geom. Dyn. {\bf 20} (2026), no. 2, 449–-474.

    \bibitem{JT24} A. Jamneshan and A. Thom, 
    Quadratic maps between non-abelian groups, 
    Math. Proc. Camb. Phil. Soc. {\bf 181} (2026), 813--837.

    \bibitem{Lac07}
    M. Lackenby,
    Adding high powered relations to large groups,
    Math. Res. Lett. {\bf 14} (2007), no. 6, 983--993. 

    \bibitem{Lei02}
    A. Leibman, 
    Polynomial mappings of groups, 
    Israel J. Math. {\bf 129} (2002), 29–-60. 

    % \bibitem{LL19}
    % A. Levit, and A. Lubotzky, 
    % Infinitely presented permutation stable groups and
    % invariant random subgroups of metabelian groups,
    % Ergodic Theory and Dynamical Systems. {\bf 42} (2022), no. 6, 2028--2063.

    % \bibitem{LL22}
    % A. Levit, and A. Lubotzky, 
    % Uncountably many permutation stable groups, 
    % Isr. J. Math. {\bf 251} (2022), 657--678.

    \bibitem{Ope26}
    OpenAI, 
    A counterexample to the soficity conjecture, 
    \url{https://openai.com/index/ten-advances-in-mathematics/}, 2026.

\end{thebibliography}
\end{document}
